\section{Possible future directions}

There are multiple future directions being pursued for \MLIR. From the more direct, reference implementation descriptions for Ops with shape functions extracted from the reference implementation, to more unknown domains, extracting equivalent transformations between Ops using the reference implementations. We've spoken about different language frontends generating \MLIR, but a missing one, and also missing from the \figref{todo} is a C++/clang mid-level representation. A, say, ``C Intermediate Language'' representation would allow easier optimization of common C++ idioms that currently need to reconstructed from lowered code (e.g., treating \lstinline{std::vector} as an array rather than pointer manipulation).

Exploring parallel and concurrency constructs in LLVM have been difficult, primarily as the changes required are invasive and not easily layered (e.g., injecting metadata and inspecting all passes to ensure metadata is propagated while losing optimization opportunities as the level is too low), this results in worse development experience for all. With \MLIR those structured constructs can be first class operations with regions with verification. This would allow the higher level transformations before lowering to, say, LLVM where regular transformations can be performed on already lowered code.

While roundtripping via textual form is useful for debugging and testing, it also useful for education. Additional tooling to show the interaction of optimizations in high-performance compilation could demystify compilers for new students.
